\subsection{Continuous operators and complete sensitivity proofs}
\label{supp:theory}

The Lipschitz results below concern the continuous prediction and allocation
operators before execution.  They exclude integer priority rounding, packet deadline switches, OpenDSS
acceptance tests and bisection, station flow active set changes, finite route
argmin selection, and integer completion events.  Finite route selection is
treated by a score error bound, while execution constraints are treated by a
separate preservation proposition.  Algebraic equalities and bounds are stated
in real arithmetic.  Electrical feasibility means satisfaction of the stated
numerical acceptance criteria.

To define the analytical map explicitly, hold the forecast continuation fixed
and set the packet and electrical factors to one.  Write
$b=e+\rho_q q$ and let $\bar s$ be the resulting charging service.  Then
\begin{equation}
F_{c,h}(x,a,\zeta)=
(\hat m_{h+1},\hat e_{h+1},b-\bar s).
\label{si:eq:continuous-service}
\end{equation}
Let $L_c$ be \eqref{eq:stage} evaluated with this same ungated service.
Neither $F_{c,h}$ nor $L_c$ includes a variable packet fraction, an AC acceptance
decision, or a crew event.  The subscript $h$ is suppressed where the same
bounds apply at every step.  Within this continuous analysis, $x$ is driven by
the fixed forecast, not by unrevealed future demand.  Forecast histories are common to the compared
trajectories.  The forecaster argument also establishes continuity when those
histories vary over a bounded set.

The analysis uses the Euclidean norm.  The projection onto a closed convex set
is nonexpansive, so
\begin{equation}
\|\Pi_{\mathcal{Z}}(u)-\Pi_{\mathcal{Z}}(v)\|\le \|u-v\|.
\label{si:eq:projection}
\end{equation}

\textit{Condition 1.}  On the stated bounded continuous prediction domain and
for all feasible requested actions, the one step service state operator
$F_c(x,a,\zeta)$ is
Lipschitz in $(x,\zeta)$ with finite constants $L_x$ and $L_z$.
\begin{equation}
\|F_c(x,a,\zeta)-F_c(\tilde x,a,\tilde \zeta)\|
\le L_x\|x-\tilde x\|+L_z\|\zeta-\tilde \zeta\|.
\end{equation}

\textit{Condition 2.}  The continuous stage risk $L_c$ is Lipschitz in
$(x,\zeta)$ with constants $C_x$ and $C_z$ for all feasible actions.
\begin{equation}
|L_c(x,a,\zeta)-L_c(\tilde x,a,\tilde\zeta)|
\le C_x\|x-\tilde x\|+C_z\|\zeta-\tilde \zeta\|.
\end{equation}
\textit{Lemma 2.}  Conditions 1 and 2 hold for the continuous service equations
on the stated bounded domain.  The continuous state transition and stage risk
are also Lipschitz in the feasible action argument.

\textit{Proof.}  Let $\mathcal D$ denote the product of the closed state,
threat, and feasible action intervals.  The backlog bound in Section III gives
a finite upper bound for $q$ because $0\le\rho_q<1$ and arrivals are bounded.
All remaining state variables have finite upper bounds by the definition of
$\mathcal D$.  The scalar maps $r\mapsto[r]_+$ and
$r\mapsto\min\{c,r\}$ are nonexpansive.  Euclidean projection and componentwise
clipping are also nonexpansive.  Every affine map has a finite operator norm.

For bounded scalars $a,a',b,b'$, the product satisfies
\begin{equation}
|ab-a'b'|\le |a|\,|b-b'|+|b'|\,|a-a'|.
\end{equation}
The right hand side has finite coefficients on $\mathcal D$.  The only quotient
in the service equations has the form $g(\ell,c)=\ell/(1+c)$.  If
$0\le\ell,\ell'\le\bar\ell$ and $c,c'\ge0$, then
\begin{align}
|g(\ell,c)-g(\ell',c')|
&\le |\ell-\ell'|
+\bar\ell |c-c'| .
\end{align}
This follows by adding and subtracting $\ell'/(1+c)$ and using
$(1+c)^{-1}\le1$ together with
$|(1+c)^{-1}-(1+c')^{-1}|\le|c-c'|$.
Explicit constants follow from these inequalities.  Use common zone bounds
$\bar m$, $\bar e$, and $\bar q$, and set $\bar b=\bar e+\rho_q\bar q$ and
$\bar c=\beta_m\bar m+\beta_e\bar b+\beta_r+\beta_p$.
For a zone input $(m,e,q,z^r,z^p,z^c,u^{\rm ch},u^{\rm com})$, valid constants are
\begin{align}
K_b&=\sqrt{1+\rho_q^2},\\
K_c&=\|(\beta_m,\beta_e,\beta_e\rho_q,\beta_r,\beta_p)\|_2,\\
K_{\ell}&=K_c+1+B_{\rm com}(1+\alpha_c),\\
K_{\eta}&=K_{\ell}+\bar cK_c,\\
K_s&=K_b+1+B_{\rm ch}(1+K_{\eta}).
\label{si:eq:service-constants}
\end{align}
The product bound gives the communication capacity constant in $K_\ell$.
Since $0\le\ell^c\le c\le\bar c$, the quotient bound gives $K_\eta$.
Applying the product bound to $u^{\rm ch}(1-z^p)\eta^c$ and then taking the
minimum with $b$ gives $K_s$.  For the remaining cost terms, valid constants are
\begin{align}
K_m&=1+\beta_{mc}+\bar m(1+\beta_{mc}K_\eta),\\
K_p&=\bar b+K_b,\qquad K_z=3+K_\eta,\\
K_L&=\lambda_mK_m+\lambda_e(K_b+K_s)+\lambda_cK_\ell
\nonumber\\
&\quad+\lambda_pK_p+\lambda_zK_z.
\label{si:eq:cost-constants}
\end{align}
These bounds follow by applying the same product inequality to mobility loss,
power loss, and the two cross domain products, respectively.
The only variable output of \eqref{si:eq:continuous-service} is $b-\bar s$.
Summing squared zone differences therefore gives the joint constant
$K_F=K_b+K_s$ for $F_c$.  Cauchy--Schwarz over the zones gives
$\sqrt N K_L$ for $L_c$.  Thus the choices
$L_x=L_z=L_a=K_F$ and $C_x=C_z=C_a=\sqrt N K_L$ are finite and satisfy
Conditions 1 and 2, including their action bounds.  \hfill $\square$

\textit{Lemma 3.}  For fixed action sequence and disturbance sequence, if two
rollout trajectories start from threat fields $\zeta_0$ and $\tilde\zeta_0$ and
from the same physical state, then
\begin{equation}
\|\zeta_h-\tilde\zeta_h\|
\le\|M\|^h\|\zeta_0-\tilde\zeta_0\|.
\label{si:eq:lemma1}
\end{equation}

\textit{Proof.}  Let $\Delta_h=\|\zeta_h-\tilde\zeta_h\|$.
Because the action and innovation sequences are
shared, \eqref{eq:propagation} and the nonexpansiveness in
\eqref{si:eq:projection} give
$\Delta_{h+1}\le\|M\|\Delta_h$.  Iteration proves \eqref{si:eq:lemma1}.
\hfill $\square$

\textit{Lemma 4.}  Under Condition 1, the service state error satisfies
\begin{equation}
E_h\le
\sum_{k=0}^{h-1}L_x^{h-1-k}L_z\Delta_k .
\label{si:eq:lemma2}
\end{equation}

\textit{Proof.}  The two rollouts use the same action sequence.  Condition 1
gives $E_{h+1}\le L_xE_h+L_z\Delta_h$.  Since the initial physical state is the
same, $E_0=0$.  Expanding the recursion gives
$E_1\le L_z\Delta_0$ and
$E_2\le L_xL_z\Delta_0+L_z\Delta_1$.  Repeating the expansion gives
\eqref{si:eq:lemma2}.  \hfill $\square$

\textit{Theorem 1.}  Let $J_H(x_0,\zeta_0)$ be the discounted $H$ step sum of $L_c$ for
a fixed feasible action sequence, and let $\tilde J_H(x_0,\tilde\zeta_0)$ be the
corresponding risk for an initial threat estimate $\tilde\zeta_0$.  For any
$\alpha>0$, define
\begin{equation}
\chi_{\alpha}=
\max\left\{\|M\|+\alpha L_z,\,
L_x\right\}.
\end{equation}
If $\|\zeta_0-\tilde\zeta_0\|\le \varepsilon$, then
\begin{equation}
|J_H(x_0,\zeta_0)-\tilde J_H(x_0,\tilde\zeta_0)|
\le
\varepsilon \left(C_z+\frac{C_x}{\alpha}\right)
\sum_{h=0}^{H-1}(\gamma\chi_{\alpha})^h .
\label{si:eq:riskbound}
\end{equation}
For every finite $H$ the sum is finite.  If $\gamma\chi_{\alpha}<1$, it is also
uniformly bounded as the horizon increases by
\begin{equation}
\varepsilon \left(C_z+\frac{C_x}{\alpha}\right)
\frac{1}{1-\gamma\chi_{\alpha}}.
\end{equation}

\textit{Proof.}  Let
$\Delta_h=\|\zeta_h-\tilde\zeta_h\|$ and
$E_h=\|x_h-\tilde x_h\|$.  Lemma 3 and Condition 1 give the one step bounds
\begin{align}
\Delta_{h+1} &\le \|M\|\Delta_h, \label{si:eq:delta_step}\\
E_{h+1} &\le L_z\Delta_h+L_xE_h . \label{si:eq:state_step}
\end{align}
For a fixed $\alpha>0$, define the weighted error
$V_h=\Delta_h+\alpha E_h$.  Multiplying \eqref{si:eq:state_step} by $\alpha$ and
adding \eqref{si:eq:delta_step} gives
\begin{align}
V_{h+1}
&\le (\|M\|+\alpha L_z)\Delta_h
+\alpha L_xE_h \nonumber\\
&\le \chi_{\alpha}\Delta_h+\alpha\chi_{\alpha}E_h
=\chi_{\alpha}V_h .
\end{align}
The two rollouts start from the same physical state, so $E_0=0$ and
$V_0=\Delta_0\le\varepsilon$.  Induction gives
$V_h\le\chi_{\alpha}^h\varepsilon$ for every $h$.  Since
$\Delta_h\le V_h$ and $E_h\le V_h/\alpha$, Condition 2 yields
\begin{equation}
|L_c(x_h,a_h,\zeta_h)-L_c(\tilde x_h,a_h,\tilde\zeta_h)|
\le
\left(C_z+\frac{C_x}{\alpha}\right)
\chi_{\alpha}^h\varepsilon .
\end{equation}
Multiplying by $\gamma^h$ and summing from $h=0$ to $H-1$ gives
\eqref{si:eq:riskbound}.  The infinite geometric series gives the uniform bound
when $\gamma\chi_{\alpha}<1$.  \hfill $\square$

\textit{Condition 3.}  The continuous service transition and stage risk are
Lipschitz in the action argument on the compact feasible set.  The PC allocation
map $\pi_c$ before execution is Lipschitz in the forecast driven service state
and threat state.  Constants are taken uniformly over the finite set of
remaining episode horizons.
There are finite constants $L_a$, $C_a$, $L_{\pi}^x$, and $L_{\pi}^{z}$ such
that
\begin{align}
\|F_c(x,a,\zeta)-F_c(x,\tilde a,\zeta)\|&\le L_a\|a-\tilde a\|,\\
|L_c(x,a,\zeta)-L_c(x,\tilde a,\zeta)|&\le C_a\|a-\tilde a\|,\\
\|\pi_c(x,\zeta,A)-\pi_c(\tilde x,\tilde\zeta,A)\|
&\le L_{\pi}^x\|x-\tilde x\|+L_{\pi}^{z}\|\zeta-\tilde\zeta\|.
\end{align}
\textit{Lemma 5.}  Condition 3 holds for the continuous PC map on the stated
compact domain.

\textit{Proof.}  Lemma 2 already establishes the first two inequalities in the
action argument.  It remains to prove the bound for $\pi_c$.  For a score vector
$v\in\mathbb R_+^N$, set $r_i(v)=\sqrt{v_i+\epsilon}$,
$S(v)=\sum_i r_i(v)$, and $p_i(v)=r_i(v)/S(v)$.  The mean value theorem gives
\begin{equation}
\|r(v)-r(w)\|_2\le\frac{1}{2\sqrt{\epsilon}}\|v-w\|_2.
\end{equation}
Moreover, $S(v)\ge N\sqrt{\epsilon}$ and
$|S(v)-S(w)|\le\sqrt N\|r(v)-r(w)\|_2$.  Since
$\|r(w)\|_2\le S(w)$, let $\Delta_S=|S(v)-S(w)|$.  Adding and subtracting
$r(w)/S(v)$ gives
\begin{align}
\|p(v)-p(w)\|_2
&\le \frac{\|r(v)-r(w)\|_2}{S(v)}
+\frac{\|r(w)\|_2\Delta_S}{S(v)S(w)} \nonumber\\
&\le\left(\frac{1}{2N\epsilon}
+\frac{1}{2\sqrt N\epsilon}\right)\|v-w\|_2.
\label{si:eq:allocator-lipschitz}
\end{align}
Equation \eqref{eq:allocator} is an affine function of $p(v)$.  Its Lipschitz
constant is therefore at most
\begin{equation}
(1-\eta_B)B
\left(\frac{1}{2N\epsilon}+\frac{1}{2\sqrt N\epsilon}\right).
\end{equation}

The forecaster requires several additional operators.  Every trained affine map
has its finite matrix norm as a Lipschitz constant.  The sigmoid and hyperbolic
tangent have derivative bounds $1/4$ and $1$.  The GRU combines these maps with
gated products.  Its hidden inputs remain bounded over the finite history, so
the product inequality in Lemma 2 applies to every gate multiplication.

For LayerNorm with hidden width $d$, write $P_d=I-\mathbf1\mathbf1^\top/d$,
$v=P_dh$, and $s=(\|v\|_2^2/d+\epsilon_{\rm LN})^{1/2}$.  Its normalization
Jacobian is
\begin{equation}
\left(s^{-1}I-\frac{vv^\top}{d s^3}\right)P_d.
\label{si:eq:layernorm-bound}
\end{equation}
At $v=0$, the bracket is $s^{-1}I$.  Otherwise, it has eigenvalues $s^{-1}$ perpendicular to $v$ and
$\epsilon_{\rm LN}/s^3$ along $v$.  Both are at most
$1/\sqrt{\epsilon_{\rm LN}}$.  With learned scale $g$ and offset $b_{\rm LN}$,
the LayerNorm constant is at most $\|g\|_\infty/\sqrt{\epsilon_{\rm LN}}$.
Its output norm is at most $\|g\|_\infty\sqrt d+\|b_{\rm LN}\|_2$, which also
provides a bound for the next recurrent input.

Exact GELU is $t\Phi(t)$, where $\Phi$ and $\phi$ are the standard normal
distribution and density.  Its derivative is $\Phi(t)+t\phi(t)$, whose absolute
value is at most $1+1/\sqrt{2\pi e}$.
Dropout is the identity during inference.  Input $\log(1+y)$ has derivative at
most one for $y\ge0$, and normalization uses fixed strictly positive scales.
The inverse transformation $\exp(\sigma o+\mu)-1$ has bounded derivative on the
bounded network output.  Final positive clipping is nonexpansive.  These facts,
the hidden state bound, and induction over the fixed history prove that the
entire implemented forecaster has a finite constant on the stated domain.

The forecast exposure weights use positive fixed training means, sums, products,
and clipping.  They are bounded and Lipschitz by Lemma 2.
For every candidate, \eqref{eq:propagation} composes an affine map, the
$1/\delta$ Lipschitz saturation, and a nonexpansive projection.
Finite iteration and the bounded product inequality therefore prove that
\eqref{eq:envelope} is Lipschitz in $(x,\zeta)$.
Taking differences and positive parts preserves this property.
For finite families $f_j$, the inequality
$|\min_j f_j(v)-\min_j f_j(w)|\le\max_j|f_j(v)-f_j(w)|$
gives the largest component constant for their minimum.
The charging, communication, central PC, and robust priority scores are
therefore Lipschitz.
Combining their constants with \eqref{si:eq:allocator-lipschitz} gives finite
$L_{\pi}^x$ and $L_{\pi}^z$.  This proof ends at the continuous allocation.
Finite route selection, packet events, the OpenDSS switch, and integer crew
completion are separate operations.  \hfill $\square$

\textit{Corollary 1 (continuous sensitivity before execution).}  Let
$J_H^{\pi_c}(x_0,\zeta_0)$ and
$\tilde J_H^{\pi_c}(x_0,\tilde\zeta_0)$ be the predictive risks generated by
the continuous PC score and allocation layer from the same physical state and
two initial threat estimates.  Define
\begin{align}
\bar L_x&=L_x+L_aL_{\pi}^x,\quad
\bar L_z=L_z+L_aL_{\pi}^{z},\\
\bar C_x&=C_x+C_aL_{\pi}^x,\quad
\bar C_z=C_z+C_aL_{\pi}^{z}.
\end{align}
For any $\alpha>0$, let
\begin{equation}
\bar\chi_{\alpha}=
\max\left\{\|M\|+\frac{\|R\|L_{\pi}^{z}}{\delta}+\alpha\bar L_z,\,
\bar L_x+\frac{\|R\|L_{\pi}^{x}}{\alpha\delta}\right\}.
\end{equation}
If $\|\zeta_0-\tilde\zeta_0\|\le\varepsilon$, then
\begin{equation}
|J_H^{\pi_c}(x_0,\zeta_0)-\tilde J_H^{\pi_c}(x_0,\tilde\zeta_0)|
\le
\varepsilon\left(\bar C_z+\frac{\bar C_x}{\alpha}\right)
\sum_{h=0}^{H-1}(\gamma\bar\chi_{\alpha})^h .
\end{equation}

\textit{Proof.}  Let $a_h=\pi_c(x_h,\zeta_h,A)$ and
$\tilde a_h=\pi_c(\tilde x_h,\tilde\zeta_h,A)$.  Condition 3 gives
\[
\|a_h-\tilde a_h\|\le L_{\pi}^xE_h+L_{\pi}^{z}\Delta_h .
\]
Combining this bound with Conditions 1 and 3 and the restoration term in
\eqref{eq:propagation} yields
\begin{align}
E_{h+1}&\le \bar L_xE_h+\bar L_z\Delta_h,\\
\Delta_{h+1}&\le
(\|M\|+\|R\|L_{\pi}^{z}/\delta)\Delta_h
+\|R\|L_{\pi}^{x}E_h/\delta .
\end{align}
For $V_h=\Delta_h+\alpha E_h$, the preceding two inequalities give
$V_{h+1}\le\bar\chi_{\alpha}V_h$.  Since $V_0\le\varepsilon$, induction gives
$V_h\le\bar\chi_{\alpha}^h\varepsilon$.  The stage risk difference satisfies
\[
|L_c(x_h,a_h,\zeta_h)-L_c(\tilde x_h,\tilde a_h,\tilde\zeta_h)|
\le \bar C_xE_h+\bar C_z\Delta_h .
\]
Using $E_h\le V_h/\alpha$ and $\Delta_h\le V_h$, then summing the discounted
stage differences, gives the stated bound.  \hfill $\square$
